\documentclass[12pt,a4paper]{article}

\input{/home/tabaka/Documents/GitHub/Defs/defs.tex}

\setcounter{zadaniaRozdzial}{3}
\setcounter{section}{0}

\begin{document}
	
	\begin{center}
		\Large Wymagania edukacyjne z matematyki w klasie 3 Liceum Ogólnokształcącego na poszczególne oceny.
	\end{center}
	
	Wymagania edukacyjne na poszczególne oceny są stopniowane.
	Uzyskanie danej oceny wymaga spełnienia wymagań przewidzianych
	dla wszystkich ocen niższych, z uwzględnieniem wymagań
	określonych dla danej oceny.
	
	\section{Funkcja wykładnicza i funkcja logarytmiczna}
	
	\begin{longtable}{|p{16cm}|}
		\hline
		
		\textbf{Wymagania edukacyjne na ocenę dopuszczającą lub dostateczną} \\ \hline
		
		\textbullet\ zapisuje daną liczbę w postaci potęgi o wykładniku wymiernym; \\
		
		\textbullet\ oblicza potęgi o wykładnikach wymiernych; \\
		
		\textbullet\ zapisuje daną liczbę w postaci potęgi o podanej podstawie i wykładniku rzeczywistym; \\
		
		\textbullet\ upraszcza wyrażenia, stosując prawa działań na potęgach, w prostych przypadkach; \\
		
		\textbullet\ oblicza wartości funkcji wykładniczej dla podanych argumentów; \\
		
		\textbullet\ sprawdza, czy dany punkt należy do wykresu funkcji wykładniczej; \\
		
		\textbullet\ wyznacza wzór funkcji wykładniczej na podstawie współrzędnych punktu należącego do jej wykresu oraz szkicuje ten wykres; \\
		
		\textbullet\ szkicuje wykres funkcji wykładniczej i podaje jej własności; \\
		
		\textbullet\ szkicuje wykres funkcji wykładniczej, przesuwając wykres odpowiedniej funkcji wzdłuż osi układu współrzędnych, oraz podaje własności otrzymanej funkcji; \\
		
		\textbullet\ oblicza logarytm danej liczby; \\
		
		\textbullet\ stosuje równości wynikające z definicji logarytmu do prostych obliczeń; \\
		
		\textbullet\ odczytuje z tablic przybliżone wartości logarytmów dziesiętnych; \\
		
		\textbullet\ stosuje twierdzenia o logarytmie iloczynu, ilorazu i potęgi do obliczania wartości wyrażeń zawierających logarytmy, w prostych przypadkach; \\
		
		\textbullet\ szkicuje wykres funkcji logarytmicznej i określa jej własności; \\
		
		\textbullet\ wyznacza wzór funkcji logarytmicznej, gdy dane są współrzędne punktu należącego do jej wykresu; \\
		
		\textbullet\ wyznacza zbiór wartości funkcji logarytmicznej określonej na podanej dziedzinie, w prostych przypadkach; \\
		
		\textbullet\ szkicuje wykres funkcji logarytmicznej, przesuwając wykres odpowiedniej funkcji wzdłuż osi układu współrzędnych; \\
		
		\textbullet\ rozwiązuje proste zadania osadzone w kontekście praktycznym, wykorzystując własności funkcji wykładniczej lub logarytmicznej. \\ \hline
		
		\textbf{Wymagania edukacyjne na ocenę dobrą lub bardzo dobrą} \\ \hline
		
		\textbullet\ upraszcza wyrażenia, stosując prawa działań na potęgach, w trudniejszych przypadkach; \\
		
		\textbullet\ porównuje liczby przedstawione w postaci potęg, korzystając z monotoniczności funkcji wykładniczej, w trudniejszych przypadkach; \\
		
		\textbullet\ odczytuje z wykresu funkcji wykładniczej zbiór rozwiązań nierówności; \\
		
		\textbullet\ wyjaśnia, jak należy przekształcić wykres funkcji, aby otrzymać wykres innej funkcji; \\
		
		\textbullet\ wyznacza podstawę logarytmu lub liczbę logarytmowaną, gdy dana jest wartość logarytmu, oraz podaje odpowiednie założenia dotyczące podstawy i liczby logarytmowanej; \\
		
		\textbullet\ stosuje twierdzenia o logarytmie iloczynu, ilorazu i potęgi do uzasadniania równości wyrażeń; \\
		
		\textbullet\ odczytuje z wykresu funkcji logarytmicznej zbiór rozwiązań nierówności; \\
		
		\textbullet\ wykorzystuje własności funkcji wykładniczej i logarytmicznej do rozwiązywania zadań osadzonych w kontekście praktycznym, np. dotyczących wzrostu wykładniczego i rozpadu promieniotwórczego; \\
		
		\textbullet\ rozwiązuje zadania dotyczące monotoniczności funkcji logarytmicznej, w tym zadania z parametrem; \\
		
		\textbullet\ dowodzi niewymierności wybranych liczb, np. $\log_2 3$. \\ \hline
		
		\textbf{Wymagania edukacyjne na ocenę celującą} \\ \hline
		
		\textbullet\ rozwiązuje zadania o znacznym stopniu trudności dotyczące funkcji wykładniczej i logarytmicznej; \\
		
		\textbullet\ przeprowadza dowody twierdzeń dotyczących działań na logarytmach. \\ \hline
		
	\end{longtable}
	
	\section{Geometria analityczna}
	
	\begin{longtable}{|p{16cm}|}
		\hline
		
		\textbf{Wymagania edukacyjne na ocenę dopuszczającą lub dostateczną} \\ \hline
		
		\textbullet\ oblicza odległość między punktami w układzie współrzędnych; \\
		
		\textbullet\ stosuje wzór na odległość między punktami do rozwiązywania zadań dotyczących wielokątów, w prostych przypadkach; \\
		
		\textbullet\ wyznacza współrzędne środka odcinka, gdy dane są współrzędne jego końców; \\
		
		\textbullet\ stosuje wzór na współrzędne środka odcinka do rozwiązywania zadań, w prostych przypadkach; \\
		
		\textbullet\ zapisuje równanie okręgu o danym środku i promieniu; \\
		
		\textbullet\ odczytuje współrzędne środka i długość promienia okręgu na podstawie kanonicznej postaci jego równania; \\
		
		\textbullet\ wyznacza równanie okręgu o danym środku, przechodzącego przez dany punkt; \\
		
		\textbullet\ sprawdza, czy dany punkt należy do danego okręgu; \\
		
		\textbullet\ rozpoznaje figury osiowosymetryczne i środkowosymetryczne; \\
		
		\textbullet\ wyznacza współrzędne obrazów punktów w symetrii osiowej względem osi układu współrzędnych oraz w symetrii środkowej względem początku układu współrzędnych. \\ \hline
		
		\textbf{Wymagania edukacyjne na ocenę dobrą lub bardzo dobrą} \\ \hline
		
		\textbullet\ stosuje wzory na odległość między punktami i współrzędne środka odcinka do rozwiązywania zadań dotyczących wielokątów, w trudniejszych przypadkach; \\
		
		\textbullet\ stosuje równanie okręgu do rozwiązywania zadań w bardziej złożonych przypadkach; \\
		
		\textbullet\ stosuje własności symetrii osiowej i środkowej do rozwiązywania zadań w trudniejszych przypadkach. \\ \hline
		
		\textbf{Wymagania edukacyjne na ocenę celującą} \\ \hline
		
		\textbullet\ rozwiązuje zadania z geometrii analitycznej o znacznym stopniu trudności. \\ \hline
		
	\end{longtable}
	
	\section{Ciągi}
	
	\begin{longtable}{|p{16cm}|}
		\hline
		
		\textbf{Wymagania edukacyjne na ocenę dopuszczającą lub dostateczną} \\ \hline
		
		\textbullet\ wyznacza kolejne wyrazy ciągu, gdy dane jest kilka jego początkowych wyrazów; \\
		
		\textbullet\ wyznacza wyrazy ciągu opisanego słownie; \\
		
		\textbullet\ szkicuje wykres ciągu; \\
		
		\textbullet\ wyznacza wzór ogólny ciągu, gdy dane jest kilka jego początkowych wyrazów; \\
		
		\textbullet\ wyznacza wskazane wyrazy ciągu określonego wzorem ogólnym; \\
		
		\textbullet\ wyznacza wyrazy ciągu spełniające dany warunek, np. przyjmujące określoną wartość, w prostych przypadkach; \\
		
		\textbullet\ podaje przykłady ciągów monotonicznych, których wyrazy spełniają podane warunki; \\
		
		\textbullet\ uzasadnia, że dany ciąg nie jest monotoniczny; \\
		
		\textbullet\ wyznacza wyraz $a_{n+1}$ ciągu określonego wzorem ogólnym; \\
		
		\textbullet\ bada monotoniczność ciągu w prostych przypadkach; \\
		
		\textbullet\ wyznacza początkowe wyrazy ciągu określonego rekurencyjnie; \\
		
		\textbullet\ wyznacza wzór rekurencyjny ciągu na podstawie jego wzoru ogólnego w prostych przypadkach; \\
		
		\textbullet\ podaje przykłady ciągów arytmetycznych; \\
		
		\textbullet\ wyznacza wyrazy ciągu arytmetycznego, gdy dane są jego pierwszy wyraz i różnica; \\
		
		\textbullet\ określa monotoniczność ciągu arytmetycznego; \\
		
		\textbullet\ wyznacza wzór ogólny ciągu arytmetycznego, gdy dane są dwa jego wyrazy; \\
		
		\textbullet\ stosuje zależność między trzema kolejnymi wyrazami ciągu arytmetycznego do wyznaczania jego wyrazów; \\
		
		\textbullet\ sprawdza, czy dany ciąg jest arytmetyczny, w prostych przypadkach; \\
		
		\textbullet\ oblicza sumę $n$ początkowych wyrazów ciągu arytmetycznego; \\
		
		\textbullet\ podaje przykłady ciągów geometrycznych; \\
		
		\textbullet\ wyznacza wyrazy ciągu geometrycznego, gdy dane są jego pierwszy wyraz i iloraz; \\
		
		\textbullet\ wyznacza wzór ogólny ciągu geometrycznego, gdy dane są dwa jego wyrazy; \\
		
		\textbullet\ określa monotoniczność ciągu geometrycznego; \\
		
		\textbullet\ sprawdza, czy dany ciąg jest geometryczny, w prostych przypadkach; \\
		
		\textbullet\ oblicza sumę $n$ początkowych wyrazów ciągu geometrycznego; \\
		
		\textbullet\ stosuje własności ciągu arytmetycznego i geometrycznego w zadaniach różnego typu, w prostych przypadkach; \\
		
		\textbullet\ oblicza wysokość kapitału przy różnych okresach kapitalizacji odsetek; \\
		
		\textbullet\ oblicza oprocentowanie lokaty i okres oszczędzania w prostych przypadkach. \\ \hline
		
		\textbf{Wymagania edukacyjne na ocenę dobrą lub bardzo dobrą} \\ \hline
		
		\textbullet\ wyznacza wzór ogólny ciągu spełniającego podane warunki w trudniejszych przypadkach; \\
		
		\textbullet\ bada monotoniczność ciągów; \\
		
		\textbullet\ wyznacza wzór rekurencyjny ciągu na podstawie jego wzoru ogólnego w trudniejszych przypadkach; \\
		
		\textbullet\ rozwiązuje zadania o podwyższonym stopniu trudności dotyczące ciągów określonych rekurencyjnie; \\
		
		\textbullet\ rozwiązuje zadania z parametrem dotyczące monotoniczności ciągu; \\
		
		\textbullet\ stosuje własności ciągu arytmetycznego oraz wzór na sumę jego początkowych wyrazów w zadaniach o podwyższonym stopniu trudności, w tym w zadaniach tekstowych; \\
		
		\textbullet\ wyznacza wartości niewiadomych, dla których wraz z danymi liczbami tworzą one ciąg arytmetyczny lub geometryczny, w prostych przypadkach; \\
		
		\textbullet\ rozwiązuje równania z zastosowaniem wzorów na sumę początkowych wyrazów ciągu arytmetycznego i geometrycznego w trudniejszych przypadkach; \\
		
		\textbullet\ stosuje zależność między trzema kolejnymi wyrazami ciągu geometrycznego w zadaniach różnego typu; \\
		
		\textbullet\ rozwiązuje zadania dotyczące lokat, w tym oblicza okres oszczędzania i oprocentowanie, oraz rozwiązuje zadania związane z kredytami; \\
		
		\textbullet\ stosuje własności ciągów arytmetycznych i geometrycznych, w tym wzory na sumę $n$ początkowych wyrazów tych ciągów, w zadaniach różnego typu, także osadzonych w kontekście praktycznym, w trudniejszych przypadkach. \\ \hline
		
		\textbf{Wymagania edukacyjne na ocenę celującą} \\ \hline
		
		\textbullet\ rozwiązuje zadania o podwyższonym stopniu trudności dotyczące ciągów, w szczególności ich monotoniczności. \\ \hline
		
	\end{longtable}
	
	\section{Statystyka}
	
	\begin{longtable}{|p{16cm}|}
		\hline
		
		\textbf{Wymagania edukacyjne na ocenę dopuszczającą lub dostateczną} \\ \hline
		
		\textbullet\ oblicza średnią arytmetyczną oraz wyznacza medianę i dominantę zestawu danych; \\
		
		\textbullet\ oblicza średnią ważoną liczb o podanych wagach. \\ \hline
		
		\textbf{Wymagania edukacyjne na ocenę dobrą lub bardzo dobrą} \\ \hline
		
		\textbullet\ oblicza średnią arytmetyczną oraz wyznacza medianę i dominantę danych przedstawionych na różne sposoby; \\
		
		\textbullet\ wykorzystuje średnią arytmetyczną, medianę, dominantę i średnią ważoną do rozwiązywania zadań w trudniejszych przypadkach; \\
		
		\textbullet\ oblicza średnią arytmetyczną oraz wyznacza medianę i dominantę danych pogrupowanych i przedstawionych na różne sposoby; \\
		
		\textbullet\ rozwiązuje zadania dotyczące statystyki w trudniejszych przypadkach. \\ \hline
		
		\textbf{Wymagania edukacyjne na ocenę celującą} \\ \hline
		
		\textbullet\ rozwiązuje zadania o znacznym stopniu trudności dotyczące statystyki. \\ \hline
		
	\end{longtable}
	
\end{document}